% Generated draft: requires mathematical and source review.
\documentclass[11pt]{article}
\usepackage[margin=25mm]{geometry}
\usepackage{fontspec}
\setmainfont{Latin Modern Roman}
\usepackage{amsmath,amssymb,mathtools,unicode-math}
\setmathfont{Latin Modern Math}
\usepackage{enumitem,needspace,xurl,hyperref}
\setlength{\parindent}{0pt}
\setlength{\parskip}{5pt}
\setlength{\emergencystretch}{3em}
\newcommand{\N}{\mathbb{N}}
\newcommand{\Z}{\mathbb{Z}}
\newcommand{\Q}{\mathbb{Q}}
\newcommand{\R}{\mathbb{R}}
\newcommand{\C}{\mathbb{C}}
\newcommand{\F}{\mathbb{F}}
\newcommand{\A}{\mathbb{A}}
\newcommand{\PP}{\mathbb{P}}
\newcommand{\E}{\mathbb{E}}
\newcommand{\eps}{\varepsilon}
\newcommand{\dd}{\,\mathrm{d}}
\begin{document}
\section*{UnsolvedMath ID: 700831}
\textbf{Source problem code:} catalogue entry
\subsection{Definitions and mathematical statement}
Let \(G\) and \(H\) be groups, each finite or countably infinite, and let \(A\) be a common subgroup, with its specified embeddings into \(G\) and \(H\) understood. Assume that \(A\neq G\) and \(A\neq H\), and that
\[
\text{for every subgroup }N\leq A,\qquad
(N\trianglelefteq G\ \text{and}\ N\trianglelefteq H)\Longrightarrow N=\{1\}.
\]
Here normality is interpreted through the specified embeddings of \(A\). Let \(G*_A H\) denote the free product with amalgamation, obtained from \(G*H\) by identifying the two images of every element of \(A\).

Put \(\mathbb N=\{1,2,3,\ldots\}\), and let \(\operatorname{Sym}(\mathbb N)\) be the group of all bijections of \(\mathbb N\). An action on \(\mathbb N\) is highly transitive if, for every integer \(k\geq1\), it is transitive on the ordered \(k\)-tuples of pairwise distinct elements of \(\mathbb N\).

\textbf{Question.} For every such choice of \(G,H,A\), does there exist an injective homomorphism
\[
\rho:G*_A H\longrightarrow\operatorname{Sym}(\mathbb N)
\]
such that, for every integer \(k\geq1\) and every two ordered tuples \((x_1,\ldots,x_k)\) and \((y_1,\ldots,y_k)\), each consisting of pairwise distinct natural numbers, there exists \(w\in G*_A H\) satisfying
\[
\rho(w)(x_i)=y_i\qquad(1\leq i\leq k)?
\]
The supplied record labels this question open. Its comment dated 2013 points to partial results, but does not specify their hypotheses or conclusions; no additional restriction or exception has been inferred from those references.
\subsection{Short English statement}
Must every amalgamated free product of two at most countable groups admit a faithful highly transitive action on \(\mathbb N\), provided the amalgamated subgroup is proper in both factors and contains no nontrivial subgroup normal in both? The open status and the reference to partial results are supplied source assessments, not independently verified conclusions.
\subsection{Sources}
\begin{itemize}
\item Local UnsolvedMath registry, supplied complete source record: permanent ID \(700831\), exact source code \texttt{catalogue entry}, title \emph{Kourovka Notebook Problem 12.11}. The supplied registry status is \texttt{open}.
\item A. M. W. Glass, proposer, 1992. \emph{The Kourovka Notebook (Unsolved Problems in Group Theory)}, original problem 12.11, 12th issue (1992); the record identifies the 21st issue edition (2026), arXiv:1401.0300v48. The supplied collection description names E. I. Khukhro and V. D. Mazurov as editors.
\item The supplied comment dated 2013 cites partial results in S. V. Gunhouse, \emph{Highly transitive representations of free products on the natural numbers}, Ph.D. thesis, Bowling Green State University, 1993.
\item The same supplied 2013 comment cites K. K. Hickin, \emph{J. London Math. Soc. (2)} \textbf{46}, no. 1 (1992), 81--91. No article title is supplied.
\item The 2013 partial-results comment and the edition information above are supplied source qualifications; neither the references nor the current mathematical status have been independently verified here.
\end{itemize}
\end{document}
